% ===================================================================
%  BẢNG Số 5 — Ngôi nhà và việc xây dựng.
%
%  Traduction, faite en 2026, en vietnamien standard d'aujourd'hui,
%  sur l'ido de texto/io. Regles de la colonne : voir 10-bang-01.tex.
%  Le tableau est un DIALOGUE et l'auteur n'y numerote aucun alinea :
%  toutes les cles portent le meme « 01 », chaque replique un rang.
%
%  LE TABLEAU DES METIERS, ET LA LANGUE Y FAIT TOUT AVEC DEUX MOTS.
%  Le vietnamien n'a pas de suffixe d'agent ; il pose devant le nom du
%  travail un CLASSIFICATEUR DE PERSONNE, et la profession se lit :
%
%    thợ  l'ouvrier de metier   thợ nề le macon, thợ mộc le menuisier,
%                               thợ rèn le forgeron, thợ lợp le
%                               couvreur, thợ sơn le peintre, thợ kính
%                               le vitrier, thợ khóa le serrurier
%    người  la personne         người làm vườn le jardinier
%    ông    le monsieur         ông kiến trúc sư l'architecte
%
%  Les dix-huit metiers de la planche se nomment donc tous de la meme
%  facon, et le lecteur voit d'un coup d'oeil que ce sont des metiers.
%  Le francais de 1926 les eparpillait — macon, charpentier, zingueur,
%  platrier, tapissier — sans qu'aucune forme les rassemble.
%
%  ET LA MAISON SE LIT DE BAS EN HAUT SANS UN SEUL EMPRUNT : móng les
%  fondations, tường les murs, cửa la porte, cửa sổ la fenetre, cầu
%  thang l'escalier, mái le toit, ngói la tuile, ống khói la cheminee.
%  Seuls entrent en francais les materiaux et outils venus par le
%  chantier colonial : xi-măng, bê-tông (non employes ici), « kẽm »
%  pour le zinc, « ê-ke » l'equerre.
%
%  LA NOTE (*) SUR « BALK-O » RESTE UNE NOTE DE LEXICOGRAPHE. Guignon
%  y propose une racine ido pour la solive, et cite six langues. Le
%  vietnamien en a une, « xà nhà », et la note le dit sans rien
%  changer a ce qu'elle discute : c'est de l'ido qu'elle parle.
% ===================================================================

%%K t05-apar-1 apar
\VUsaut{6.0mm}
\VUtitre{15.0pt}{60}{BẢNG Số 5}
\VUsaut{1.4mm}
\VUtitre{11.4pt}{0}{Ngôi nhà và việc xây dựng.}
\VUsaut{2.2mm}

%%K t05-tit-1 sub
\VUpk{10.2pt}{40}{Cuộc gặp may mắn.}
\VUsaut{3.0mm}

%%K t05-01-1 p
\textsc{Gioan}. --- Bác ơi, cháu rất muốn biết người ta xây một \VUgras{ngôi nhà} như
thế nào.

%%K t05-01-2 p
\textsc{Người bác} \textsuperscript{(80)}. --- Dễ thôi, cháu ạ; đi với bác, bác sẽ chỉ
cho cháu ngôi nhà mà ông Bênađô \textsuperscript{(79)} đang cho xây và bác sẽ ở.

%%K t05-01-3 p
\textsc{Gioan}. --- Thế ai đã vẽ \VUgras{bản thiết kế} \textsuperscript{(76)} ngôi nhà
này ạ?

%%K t05-01-4 p
\textsc{Người bác}. --- \VUgras{Ông kiến trúc sư} \textsuperscript{(75)}; ông đã trình
một bản vẽ rất rõ ràng, bản vẽ được duyệt, và \VUgras{nhà thầu}
\textsuperscript{(77)} đã bắt đầu công việc.

%%K t05-01-5 p
\textsc{Gioan}. --- Thế nhà thầu là ai ạ?

%%K t05-01-6 p
\textsc{Người bác}. --- Là ông Blăng, cha của bạn Giacôbê của cháu. Ông chỉ huy thợ
thuyền cùng với ông Ru, kiến trúc sư.

%%K t05-01-7 p
Kìa! Vừa đúng lúc ông Blăng đến với \VUgras{cái thước} \textsuperscript{(78)} của mình,
và ông Ru với bản thiết kế.

%%K t05-01-8 p
Chào các ông. Các ông có khỏe không?

%%K t05-01-9 p
\textsc{Ông Ru}. --- Rất khỏe, cảm ơn ông. Chào cháu Gioan, cháu bé này lớn quá!

%%K t05-01-10 p
\textsc{Người bác}. --- Vâng, thật vậy, cháu đã ra dáng người lớn. Cháu rất nghiêm túc,
phải giảng cho cháu mọi thứ cháu trông thấy. Hôm nay cháu muốn biết người ta xây một
ngôi nhà như thế nào.

%%K t05-tit-2 sub
\VUpk{10.2pt}{40}{Thợ thuyền.}
\VUsaut{3.0mm}

%%K t05-01-11 p
\textsc{Ông Blăng}. --- Nào, đi với bác, cháu bé, bác sẽ giảng ngay cho cháu, vì bác rất
quý những đứa trẻ ham học.

%%K t05-01-12 p
Cháu thấy người thợ cầm cái \VUgras{xẻng} \textsuperscript{(82)} trên tay kia không :
đó là \VUgras{thợ đào đất} \textsuperscript{(81)}. Ông ấy đào một cái \VUgras{rãnh}
\textsuperscript{(46)} để đặt ống nước. Ông cũng đã đào đất ở chỗ đặt ngôi nhà, để thợ
nề dựng được bốn bức \VUgras{tường}. Phần tường nằm dưới đất gọi là \VUgras{móng}
\textsuperscript{(2)}. Khoảng trống giữa các móng làm thành \VUgras{tầng hầm}
\textsuperscript{(3)}. Tầng hầm ấy có một ô cửa nhỏ gọi là \VUgras{lỗ thông hơi}
\textsuperscript{(5)}, một \VUgras{cửa} \textsuperscript{(4)} và một cầu thang.

%%K t05-01-13 p
\VUgras{Thợ nề} \textsuperscript{(108)} đã xây tiếp các bức tường, chừa lại những ô
trống cho \VUgras{cửa ra vào} \textsuperscript{(8)} và \VUgras{cửa sổ}
\textsuperscript{(17)}.

%%K t05-01-14 p
\textsc{Gioan}. --- Nhưng cháu không thấy người thợ nề ấy!

%%K t05-01-15 p
\textsc{Ông Blăng}. --- Ông ấy làm xong ở ngôi nhà rồi. Ông đang ở gần \VUgras{chuồng
ngựa} \textsuperscript{(54)}, trên \VUgras{giàn giáo} \textsuperscript{(110)} của mình;
ông xây bức tường của \VUgras{nhà giặt} \textsuperscript{(56)}.

%%K t05-01-16 p
Ông có cái bay, cái máng và một \VUgras{dây dọi} \textsuperscript{(109)}. Nhưng cháu sẽ
không nhớ nổi những cái tên ấy đâu.

%%K t05-01-17 p
\textsc{Gioan}. --- Có chứ ạ, vì cháu sẽ xin bác cháu một cái bay nhỏ và một cái máng
ngay, còn dây dọi thì cháu sẽ tự làm lấy bằng sợi chỉ.

%%K t05-tit-3 sub
\VUsaut{3.0mm}
\VUpk{10.2pt}{40}{Vật liệu.}
\VUsaut{3.0mm}

%%K t05-01-18 p
\textsc{Người bác}. --- À! rất tốt. Nhưng cháu nói bác nghe, Gioan, muốn xây một ngôi
nhà thì cần những gì?

%%K t05-01-19 p
\textsc{Gioan}. --- Cần \VUgras{đá xây} \textsuperscript{(59)} và \VUgras{vữa}
\textsuperscript{(65)} ạ.

%%K t05-01-20 p
\textsc{Người bác}. --- Đúng vậy. Cháu nhìn người đàn ông đội mũ rộng vành trên chiếc
\VUgras{xe bò} \textsuperscript{(136)} kia. Đó là \VUgras{người đánh xe}
\textsuperscript{(135)}. Ngày nào ông cũng chở đến đây những \VUgras{khối đá}
\textsuperscript{(60)}. Ông dỡ xe xuống sân, rồi \VUgras{thợ đẽo đá}
\textsuperscript{(83)} đẽo các khối đá và cưa chúng bằng \VUgras{cưa đá}
\textsuperscript{(85)}. Một \VUgras{người phụ việc} \textsuperscript{(146)} mang đá đến
cho thợ nề đặt vào chỗ.

%%K t05-01-21 p
Trong lúc đó, \VUgras{người trộn vữa} \textsuperscript{(87)} chuẩn bị vữa. Ông cần
\VUgras{cát} \textsuperscript{(62)}, \VUgras{vôi} \textsuperscript{(63)} và
\VUgras{nước} \textsuperscript{(64)}. Vôi ở cạnh ông trong một cái \VUgras{bao}
\textsuperscript{(71)}; một \VUgras{thợ phụ nề} \textsuperscript{(111)} mang cát đến cho
ông bằng \VUgras{xe cút kít} \textsuperscript{(112)}, còn \VUgras{người học việc}
\textsuperscript{(113)} thì đứng ở máy bơm \textsuperscript{(58)}. Anh ta đổ đầy một
\VUgras{cái xô} \textsuperscript{(114)}. Vữa sắp xong. \VUgras{Người khiêng vữa}
\textsuperscript{(107)} lấy vữa và theo thang mang lên giàn giáo, nơi một người thợ nề
đang hoàn thành mặt này của ngôi nhà.

%%K t05-01-22 p
Còn bây giờ, người thợ làm \VUgras{mái nhà} \textsuperscript{(31)} gọi là gì?

%%K t05-01-23 p
\textsc{Gioan}. --- Là \VUgras{thợ lợp mái} \textsuperscript{(118)} ạ.

%%K t05-tit-4 sub
\VUsaut{3.0mm}
\VUpk{10.2pt}{40}{Mái nhà.}
\VUsaut{3.0mm}

%%K t05-01-24 p
\textsc{Ông Blăng}. --- Đúng, nhưng trước hết phải có \VUgras{thợ mộc kèo cột}
\textsuperscript{(115)}.

%%K t05-01-25 p
Thợ mộc kèo cột làm \VUgras{bộ khung mái} \textsuperscript{(35)}. Ông ghép các
\VUgras{xà} \textsuperscript{(37)} bằng \VUgras{chốt gỗ} \textsuperscript{(117)}. Những
người thợ trên cao kia, mặc quần nhung rộng, là thợ mộc kèo cột. Một người cầm một
thanh xà nhỏ, người kia đẽo một cái chốt bằng \VUgras{rìu} \textsuperscript{(116)}.
Người ở gần \VUgras{ống khói} \textsuperscript{(41)} là thợ lợp mái. Ông đóng những
thanh lati lên (*) xà nhà, rồi đóng \VUgras{ngói đá} \textsuperscript{(66)} lên lati.
Đôi khi ông lợp ngói đất.

%%K t05-noto-1 noto
\VUnotes{143.2mm}{%
(*) \VUgras{Balk-o} = tiếng Đức \textit{Balken}, tiếng Anh \textit{joist}, tiếng Pháp
\textit{solive}, tiếng Ý \textit{travicello}, tiếng Nga \textit{balka}, tiếng Tây Ban Nha
và tiếng Bồ Đào Nha \textit{viga}.
Một gốc từ kỹ thuật đến nay chưa được nhận, nhưng nên đề nghị để dịch nghĩa của từ Pháp
\textit{solive}, vốn không phải là xà cái (Pháp : \textit{poutre}) cũng không phải là
thanh xà nhỏ. Đó là một thanh gỗ vuông cạnh đỡ sàn nhà và trần nhà.%
}

%%K t05-01-26 p
Mái chuồng ngựa lợp \VUgras{ngói đất} \textsuperscript{(55)}, mái ngôi nhà lợp
\VUgras{ngói đá} \textsuperscript{(32)} và mái hiên lợp \VUgras{kẽm}
\textsuperscript{(24)}.

%%K t05-01-27 p
\textsc{Gioan}. --- Thợ lợp mái có làm cả mái kẽm không ạ?

%%K t05-01-28 p
\textsc{Ông Blăng}. --- Không, đó là \VUgras{thợ kẽm} \textsuperscript{(121)} hay
\VUgras{thợ ống nước}, người thợ đang lắp một \VUgras{cửa sổ mái}
\textsuperscript{(38)} trên mái nhà kia. Ông cũng lắp \VUgras{máng xối}
\textsuperscript{(43)} và \VUgras{ống máng} \textsuperscript{(42)}. Ông hàn kẽm với sắt
tây. Chính ông đã gắn \VUgras{cột thu lôi} \textsuperscript{(40)} và \VUgras{chong
chóng gió}\textsuperscript{(39)} lên \VUgras{đầu hồi} \textsuperscript{(36)}.

%%K t05-01-29 p
\textsc{Gioan}. --- Vậy là ngôi nhà xong rồi ạ?

%%K t05-tit-5 sub
\VUsaut{3.0mm}
\VUpk{10.2pt}{40}{Đồ nghề.}
\VUsaut{3.0mm}

%%K t05-01-30 p
\textsc{Ông Blăng}. --- Chưa hẳn. \VUgras{Thợ trát} \textsuperscript{(99)} xây bằng
\VUgras{gạch} \textsuperscript{(61)} những vách ngăn chia ngôi nhà thành các phòng. Ông
trát \VUgras{vữa thạch cao} \textsuperscript{(70)} lên \VUgras{trần}
\textsuperscript{(26)} và tường. Lúc này ông đang làm trần \VUgras{hiên nhà}
\textsuperscript{(33)}. Cũng như thợ nề, ông có một cái \VUgras{bay}
\textsuperscript{(100)} và một cái \VUgras{máng} \textsuperscript{(101)}.

%%K t05-01-31 p
Sau đó, thợ mộc làm cửa ra vào, cửa sổ, \VUgras{sàn gỗ} \textsuperscript{(16)} và
\VUgras{cửa chớp} \textsuperscript{(19)}.

%%K t05-01-32 p
Lúc này ông đang làm việc ngoài sân, trên \VUgras{bàn thợ} \textsuperscript{(90)} của
mình. Ông có một cái \VUgras{cưa} \textsuperscript{(94)} để cưa \VUgras{gỗ}
\textsuperscript{(69)}, một cái \VUgras{bào} \textsuperscript{(91)}, một cái \VUgras{kẹp
bàn} \textsuperscript{(92)}, một cái \VUgras{đục} \textsuperscript{(94 \textit{bis})},
một cái \VUgras{com-pa} \textsuperscript{(95)} và một cái \VUgras{vồ gỗ}
\textsuperscript{(95 \textit{bis})}.

%%K t05-01-33 p
\textsc{Gioan}. --- A! nhưng làm mộc vui hơn làm nề. Bác ơi, bác mua cho cháu một cái
bàn thợ, một cái cưa và một cái bào nhé, vì sắp tới cháu sẽ đóng một cái chuồng cho
Mêđo.

%%K t05-01-34 p
\textsc{Người bác}. --- Được, cháu ạ.

%%K t05-01-35 p
\textsc{Gioan}. --- Cháu vui quá! Thế còn người đứng gần cửa chính là ai ạ?

%%K t05-tit-6 sub
\VUsaut{3.0mm}
\VUpk{10.2pt}{40}{Việc sơn.}
\VUsaut{3.0mm}

%%K t05-01-36 p
\textsc{Ông Blăng}. --- Đó là \VUgras{thợ sơn} \textsuperscript{(102)}. Ông đứng trên
thang với hộp \VUgras{sơn} \textsuperscript{(103)} và cái \VUgras{chổi sơn}
\textsuperscript{(104)} của mình. Ông sơn gỗ cửa chính để gỗ bền lâu hơn.

%%K t05-01-37 p
Chính ông cũng là người dán giấy tường trong các căn hộ.

%%K t05-01-38 p
\VUgras{Thợ bọc} \textsuperscript{(107)}, người đang đứng trên \VUgras{thang}
\textsuperscript{(106)}, ở \VUgras{ban công} \textsuperscript{(28)}, lắp một tấm
\VUgras{mành} \textsuperscript{(29)}.

%%K t05-01-39 p
Người thợ đứng gần ô cửa sổ lớn là \VUgras{thợ kính} \textsuperscript{(96)}. Ông gắn các
\VUgras{tấm kính} \textsuperscript{(18)} bằng \VUgras{ma-tít} \textsuperscript{(73)}.
Người thợ kia, đang quỳ gần cửa tầng hầm, là \VUgras{thợ khóa}
\textsuperscript{(97)}. Ông lắp một cái \VUgras{ổ khóa} \textsuperscript{(6)}. Cái
\VUgras{giũa} \textsuperscript{(98)} của ông ở cạnh ông.

%%K t05-01-40 p
\textsc{Gioan}. --- Người thợ đang gõ \VUgras{búa} \textsuperscript{(84)} lên sắt kia
cũng là thợ khóa ạ?

%%K t05-01-41 p
\textsc{Ông Blăng}. --- Đúng hơn thì đó là thợ rèn \textsuperscript{(125)}. Ông rèn
những \VUgras{món đồ sắt} \textsuperscript{(67)} cho \VUgras{thợ điện}
\textsuperscript{(124)}, người đang ở trên mái \VUgras{nhà xe} \textsuperscript{(51)}.
Ông có một cái \VUgras{lò rèn} \textsuperscript{(126)}, một cái \VUgras{đe}
\textsuperscript{(127)} và một cái \VUgras{kìm} \textsuperscript{(128)}.

%%K t05-01-42 p
Thợ điện lắp \VUgras{dây đồng} \textsuperscript{(44)} của điện thoại.

%%K t05-tit-7 sub
\VUpk{10.2pt}{40}{Việc làm vườn.}
\VUsaut{3.0mm}

%%K t05-01-43 p
\textsc{Người bác}. --- Cuối \VUgras{sân} \textsuperscript{(45)}, gần \VUgras{hàng rào
sắt} \textsuperscript{(47)} và \VUgras{cổng xe} \textsuperscript{(48)}, cháu thấy
\VUgras{người làm vườn} \textsuperscript{(129)}. Ông đang sửa soạn \VUgras{mảnh vườn}
\textsuperscript{(49)} nhỏ. Ông đã cuốc đất bằng cái \VUgras{mai}
\textsuperscript{(131)}, ông gieo hạt và cào bằng cái \VUgras{cào}
\textsuperscript{(130)}. Sau đó, với cái \VUgras{bình tưới} \textsuperscript{(132)}, ông
sẽ tưới các \VUgras{luống hoa} \textsuperscript{(150)} và cây sẽ lớn lên.

%%K t05-01-44 p
\VUgras{Người coi ngựa} \textsuperscript{(133)}, vừa chăm sóc và cho ngựa ăn xong, lau
chiếc \VUgras{xe} \textsuperscript{(52)} bằng bọt biển, \VUgras{bơm tay}
\textsuperscript{(134)} và một xô nước. Rồi ông sẽ đưa xe vào nhà xe và đóng cửa bằng
cái \VUgras{then} \textsuperscript{(53)} to.

%%K t05-01-45 p
Bác nghĩ cháu hài lòng rồi đấy; cháu đã được giảng đủ.

%%K t05-01-46 p
\textsc{Gioan}. --- Vâng, thưa bác, nhưng cháu rất muốn xem bên trong ngôi nhà.

%%K t05-01-47 p
\textsc{Người bác}. --- Chuyện ấy để mai. Ông Ru sẽ giảng cho cháu bản thiết kế và cho
cháu biết tên từng căn phòng.

%%K t05-tit-8 sub
\VUsaut{3.0mm}
\VUpk{10.2pt}{40}{Cách bố trí bên trong ngôi nhà.}
\VUsaut{2.4mm}

%%K t05-apar-2 apar
{\centering\textit{(Xem bản thiết kế.)}\par}
\VUsaut{2.4mm}

%%K t05-01-48 p
\textsc{Ông kiến trúc sư}. --- Lại đây, Gioan, bác sẽ đưa cháu đi xem các phần khác nhau
của ngôi nhà.

%%K t05-01-49 p
Ta vào \VUgras{tầng trệt} \textsuperscript{(7)}. Kia là nút \VUgras{chuông cửa}
\textsuperscript{(11)}. Ta bước lên ba \VUgras{bậc} \textsuperscript{(14)} của
\VUgras{thềm} \textsuperscript{(9)}, ta qua \VUgras{ngưỡng} \textsuperscript{(10)} của
\VUgras{cửa chính} \textsuperscript{(8)}, và đây, ta đang ở chân \VUgras{cầu thang}
\textsuperscript{(13)}.

%%K t05-01-50 p
\textsc{Gioan}. --- Đây là hành lang ạ.

%%K t05-01-51 p
\textsc{Ông kiến trúc sư}. --- Không, đây là \VUgras{tiền sảnh} \textsuperscript{(12)}.
\VUgras{Hành lang} \textsuperscript{(a)} ở phía cuối. Bên phải là \VUgras{phòng khách}
\textsuperscript{(b)} và \VUgras{phòng ăn} \textsuperscript{(c)}; đối diện phòng ăn là
\VUgras{bếp} \textsuperscript{(d)} và \VUgras{phòng để thức ăn} \textsuperscript{(e)},
rồi phía trước là \VUgras{phòng làm việc} \textsuperscript{(f)}. \VUgras{Hiên nhà}
\textsuperscript{(23)}, với \VUgras{cửa kính} \textsuperscript{(25)}, ở cạnh phòng khách.

%%K t05-01-52 p
Bây giờ ta lên cầu thang. Cháu vịn vào \VUgras{tay vịn} \textsuperscript{(15)} và đếm
bậc nhé. Có mười bốn bậc. Đây là \VUgras{chiếu nghỉ} \textsuperscript{(m)} : ta đang ở
\VUgras{tầng một} \textsuperscript{(27)}.

%%K t05-tit-9 sub
\VUsaut{3.0mm}
\VUpk{10.2pt}{40}{Tầng một.}
\VUsaut{3.0mm}

%%K t05-01-53 p
Đối diện cầu thang là \VUgras{nhà vệ sinh} \textsuperscript{(20)} và \VUgras{phòng rửa
mặt} \textsuperscript{(j)}. Bên phải, một \VUgras{phòng ngủ} \textsuperscript{(g)} đẹp
có hai cửa sổ nhìn ra \VUgras{mặt tiền} \textsuperscript{(1)} ngôi nhà. Bên trái là
\VUgras{phòng tắm} \textsuperscript{(1)} và \VUgras{phòng dành cho khách}
\textsuperscript{(i)} có một \VUgras{hốc tường} \textsuperscript{(h)} lớn. Cạnh phòng
ngủ là \VUgras{phòng của trẻ con} \textsuperscript{(k)}. Phòng này nhìn ra một
\VUgras{sân thượng} \textsuperscript{(21)} rộng. Dưới sân thượng ấy có một nhà vệ sinh
khác \textsuperscript{(20)}, và trên tường, kia là cái \VUgras{chuông}
\textsuperscript{(22)} báo giờ ăn tối.

%%K t05-01-54 p
\textsc{Gioan}. --- Ta lên \VUgras{tầng hai} đi ạ.

%%K t05-01-55 p
\textsc{Ông kiến trúc sư}. --- Không có tầng hai. Dưới mái chỉ có \VUgras{phòng áp mái}
\textsuperscript{(33)} và kho \textsuperscript{(34)}. Ta xem xong rồi, có thể xuống được.

%%K t05-01-56 p
\textsc{Gioan}. --- Cảm ơn bác, thật thú vị. Cháu sẽ kể ngay cho cha cháu tất cả những
gì cháu đã thấy.
\VUsaut{4.0mm}
\VUfilet{22mm}
